\section{Conclusion}

HyperbolicMomentum studies how to optimize hierarchical representations when curvature itself is a learnable, group-dependent state. The historical exploratory result is that grouped learnable curvature improves selected unfiltered tail-ranking signals under a Euclidean Adam protocol, with WN18RR showing Hits@10 $+4.1\%$ and loss $-2.8\%$, FB15k-237 showing MRR $+2.2\%$ and Hits@1 $+3.1\%$, and synthetic hierarchies showing faster convergence rather than a better final optimum. The mechanism is curvature differentiation: WN18RR spreads from $K=2.000$ to $K=1.008$, FB15k-237 has a much smaller $\Kstd=0.071$, and synthetic trees order curvature by branching heterogeneity. Establishing a real-data advantage for cross-curvature momentum still requires completed controlled comparisons in architectures with separate home manifolds, followed by multi-seed and higher-precision replication.

The completed 2000-update mechanism matrix supports a conditional advantage against matched, validation-selected fixed-curvature controls: $2.94\%$ lower held-out distortion in the stationary hierarchy, and $3.31\%$ and $5.02\%$ under $2\times$ and $8\times$ smooth drift. The complete six-case matrix also contains ties and negative results. The point-only transport ablation remains close, while separate operator experiments verify the consistency gained by transporting momentum. These findings distinguish task-dependent curvature adaptation from the geometric correctness of optimizer-state transport.

The completed graph comparisons show small mixed differences and do not establish a robust additional task-family win. The amended execution budget is complete: 705 full-length configurations, with all delayed test evaluations released; 188 other registrations remain explicitly deferred. The completed real-data extension and distance-model block do not establish an additional family-level advantage. The full outcome matrix preserves these results alongside the positive conditional mechanism findings.
